\documentclass[12pt]{article}
\usepackage[utf8]{inputenc}
\usepackage[T1]{fontenc}
\usepackage{hyperref}
\title{Kunduvadi Language Resources for K-12 Curriculum Development}
\author{Compiled by Claude Sonnet 5.5 (claude-sonnet-5-5) for Andr\'as Kornai}
\date{}
\begin{document}
\maketitle

\section{Introduction}
Kunduvadi (ISO 639-3 \texttt{wku}, Glottocode kund1253) is a Southern Dravidian language of India. Wikipedia cites about 1,000 native speakers (2004). Its structure resembles Malayalam, and its closest relative is Pathiya. The search found very little about it.

\section{Primary Educational Resources}
No primers, textbooks or school programmes were found.

\section{Linguistic and Reference Materials}
\begin{itemize}
\item Wikipedia, Kunduvadi language: \url{https://en.wikipedia.org/wiki/Kunduvadi_language}
\item Glottolog languoid page (the URL seen was for id kani1275, which may not be the Kunduvadi entry): \url{https://glottolog.org/resource/languoid/id/kani1275}
\item Wikipedia, Kalanadi language (related): \url{https://en.wikipedia.org/wiki/Kalanadi_language}
\end{itemize}
No dictionary or grammar was found.

\section{Religious and Community Resources}
None found.

\section{Digital and Audio Resources}
None found.

\section{Curriculum Development Context}
The language is not known to be taught in schools. Any curriculum would have to start with basic documentation, and could draw on Malayalam and other related languages for comparison.
\end{document}